% --- Лекция 2: Продолжение исследования непрерывных цепей Маркова ---

\section{Вероятностный смысл элементов матрицы интенсивностей и механика скачков}

На предыдущей лекции мы установили смысл диагональных элементов матрицы интенсивностей $Q$: величина $q_{ii} = -q_i$ определяет интенсивность выхода из состояния $i$, а время пребывания в этом состоянии распределено экспоненциально с параметром $q_i$. Теперь перейдем к рассмотрению внедиагональных элементов $q_{ij}$ при $i \neq j$ и выясним, как они определяют вероятности переходов (скачков) между различными состояниями.

\begin{definition}{Вероятность перехода в момент скачка}
    Обозначим через $F_{ij}(h)$ условную вероятность того, что цепь Маркова в момент времени $t+h$ окажется в состоянии $j$, при условии, что в момент $t$ она находилась в состоянии $i$, и к моменту $t+h$ произошло изменение состояния (цепь не находится в $i$):
    \begin{equation}
        F_{ij}(h) \doteq \mathbb{P}(X(t+h) = j \mid X(t) = i, X(t+h) \neq i), \quad i \neq j.
    \end{equation}
\end{definition}

Рассмотрим предельную вероятность такого перехода при стремлении временного интервала $h$ к нулю.

\begin{statement}
    Для непрерывной цепи Маркова предел условной вероятности перехода в состояние $j$ в момент скачка из состояния $i$ равен отношению соответствующей интенсивности к суммарной интенсивности выхода из $i$:
    \begin{equation}
        \lim_{h \to 0} F_{ij}(h) = \frac{q_{ij}}{q_i}.
    \end{equation}
\end{statement}

\begin{sproof}
    Воспользуемся определением условной вероятности $\mathbb{P}(A \mid B) = \frac{\mathbb{P}(A \cap B)}{\mathbb{P}(B)}$. 
    Заметим, что событие $\{X(t+h) = j\}$ при $j \neq i$ влечет за собой событие $\{X(t+h) \neq i\}$. Следовательно, их пересечение совпадает с первым событием: $\{X(t+h) = j\} \cap \{X(t+h) \neq i\} = \{X(t+h) = j\}$.
    
    Тогда искомая вероятность записывается как:
    \begin{equation}
        F_{ij}(h) = \frac{\mathbb{P}(X(t+h) = j \mid X(t) = i)}{\mathbb{P}(X(t+h) \neq i \mid X(t) = i)}.
    \end{equation}
    
    Применим асимптотические разложения переходных вероятностей при $h \to 0$:
    \begin{enumerate}
        \item В числителе: $P_{ij}(h) = q_{ij} h + o(h)$ для всех $i \neq j$.
        \item В знаменателе: вероятность выхода из состояния $i$ есть дополнение к вероятности остаться в нем:
        $1 - P_{ii}(h) = 1 - (1 + q_{ii} h + o(h)) = -q_{ii} h - o(h) = q_i h + o(h)$.
    \end{enumerate}
    
    Подставляя эти разложения, получаем:
    \begin{equation}
        F_{ij}(h) = \frac{q_{ij} h + o(h)}{q_i h + o(h)} = \frac{q_{ij} + \frac{o(h)}{h}}{q_i + \frac{o(h)}{h}}.
    \end{equation}
    
    Переходя к пределу при $h \to 0$, все члены порядка $\frac{o(h)}{h}$ исчезают, и мы получаем окончательный результат:
    \begin{equation}
        \lim_{h \to 0} F_{ij}(h) = \frac{q_{ij}}{q_i}.
    \end{equation}
\end{sproof}

Таким образом, внедиагональные элементы $q_{ij}$ определяют «приоритетность» выбора конкретного состояния $j$ в момент, когда процесс покидает состояние $i$.

\begin{note}
    Данный результат согласуется с конструкцией, предложенной Кай Лаем. Можно строго доказать, что времена пребывания в состояниях независимы в совокупности, а выбор следующего состояния не зависит от времени, проведенного в текущем.
\end{note}

\subsection*{Механика поведения цепи}

На основе полученных свойств можно описать эволюцию цепи Маркова как дискретно-событийный процесс:
\begin{enumerate}
    \item Пусть в начальный момент времени $t=0$ процесс находится в некотором состоянии $i \in S$.
    \item Процесс «задерживается» в состоянии $i$ в течение случайного времени $\tau_i$, имеющего экспоненциальное распределение с параметром $q_i$ (то есть $\tau_i \sim \text{Exp}(q_i)$).
    \item В момент завершения времени ожидания процесс совершает скачок в новое состояние $j \neq i$ с вероятностью $\frac{q_{ij}}{q_i}$.
    \item После перехода в состояние $j$ процесс снова «задерживается» в нем на время $\tau_j \sim \text{Exp}(q_j)$, и цикл повторяется.
\end{enumerate}

\section{Эргодические непрерывные цепи Маркова}

Перейдем к изучению асимптотического поведения непрерывных цепей Маркова при $t \to \infty$. Как и в дискретном случае, центральное место здесь занимает понятие стационарного распределения и классификация состояний.

\begin{definition}{Стационарное распределение}
    Вектор вероятностей $\pi_0$ называется стационарным распределением непрерывной цепи Маркова, если для любого момента времени $t > 0$ выполняется уравнение:
    \begin{equation}
        P^T(t) \cdot \pi_0 = \pi_0.
    \end{equation}
    Для однородной цепи с матрицей интенсивностей $Q$ стационарное распределение находится из \textbf{стационарного уравнения Колмогорова}:
    \begin{equation}
        \boxed{Q^T \pi_0 = 0}.
    \end{equation}
\end{definition}

\subsection{Классификация состояний и их свойства}

Классификация состояний в непрерывном времени $t \in [0, +\infty)$ строится аналогично дискретному случаю, однако вместо числа шагов $n$ мы рассматриваем интегральные характеристики на всей временной оси.

\begin{definition}{Сообщаемость и существенность}
    \begin{enumerate}
        \item Состояние $j$ \textit{достижимо} из $i$ ($i \to j$), если $\exists t > 0 : p_{ij}(t) > 0$.
        \item Состояния $i$ и $j$ \textit{сообщаются} ($i \leftrightarrow j$), если $i \to j$ и $j \to i$.
        \item Состояние $i$ называется \textit{существенным}, если из того, что $i \to j$, следует $j \to i$.
    \end{enumerate}
\end{definition}

\begin{theorem}{Теорема о солидарности}
    Для неразложимой цепи Маркова (где все состояния сообщаются, $i \leftrightarrow j$) все состояния имеют один и тот же тип: они либо все невозвратны, либо все нулевые возвратные, либо все положительно возвратные.
\end{theorem}

\begin{note}
В отличие от дискретных цепей, в непрерывном случае понятие \textbf{периодичности} не вводится. Поскольку время перехода между состояниями является случайной величиной с экспоненциальным распределением, возможен возврат за любое время $t > 0$, что делает невозможным существование фиксированного шага (периода) $d > 1$. Таким образом, все состояния в непрерывных цепях апериодичны.
\end{note}

Для введения понятий возвратности обозначим через $\tau_{ii} \doteq \inf \{t > J_1 : X(t) = i \mid X(0) = i\}$ случайную величину времени первого возвращения в состояние $i$ (после первого скачка $J_1$).

\begin{definition}{Типы возвратности}
    \begin{enumerate}
        \item Состояние $i$ называется \textit{возвратным}, если $\mathbb{P}(\tau_{ii} < \infty) = 1$. Аналогично можно показать, что данное определение в терминах переходных вероятностей это эквивалентно:
            \begin{equation}
                \int_0^{+\infty} p_{ii}(t) \, dt = \infty.
            \end{equation}
        \item Если $\mathbb{P}(\tau_{ii} < \infty) < 1$ (или интеграл выше сходится), состояние называется \textit{невозвратным}.
        \item Возвратное состояние $i$ называется \textit{положительно возвратным}, если среднее время возврата конечно ($\mathbb{E}\tau_{ii} < \infty$). Это эквивалентно условию, что предельная вероятность пребывания в состоянии строго положительна:
            \begin{equation}
                \lim_{t \to \infty} p_{ii}(t) = \pi_i > 0.
            \end{equation}
        \item Возвратное состояние $i$ называется \textit{нулевым возвратным}, если среднее время возврата бесконечно ($\mathbb{E}\tau_{ii} = \infty$). В этом случае, несмотря на неизбежность возвращения, вероятность застать её в состоянии $i$ стремится к нулю:
            \begin{equation}
                \lim_{t \to \infty} p_{ii}(t) = 0.
            \end{equation}
    \end{enumerate}
\end{definition}

\begin{statement}[Классификация состояний в конечных цепях]
    В конечной непрерывной цепи Маркова:
    \begin{itemize}
        \item Все состояния в закрытых классах являются положительно возвратными и существенными.
        \item Все состояния в открытых классах являются нулевыми, невозвратными и несущественными.
        \item Нулевых возвратных состояний в конечных цепях не существует.
    \end{itemize}
\end{statement}

\subsection{Эргодические теоремы}

\begin{definition}{Сильная эргодичность}
    Непрерывная цепь Маркова называется сильно эргодической, если $\forall i, j \in S$ существует предел $\lim_{t \to \infty} p_{ij}(t) = p_j > 0$, не зависящий от начального состояния $i$.
\end{definition}

\begin{theorem}{Эргодическая теорема для конечных цепей}
    Конечная непрерывная цепь Маркова сильно эргодична тогда и только тогда, когда она неразложима.
\end{theorem}

\begin{tproof}
    $\Rightarrow$ Пусть цепь сильно эргодична. Тогда $\lim_{t \to \infty} p_{ij}(t) = p_j > 0$, что влечет существование такого $T$, что при всех $t \geqslant T$ выполняется $p_{ij}(t) > 0$. Это означает, что любое состояние достижимо из любого другого, то есть цепь неразложима.
    
    $\Leftarrow$ В дискретном случае требовалась апериодичность. Как уже было отмечено, в непрерывном времени $p_{ii}(t) > 0$ для всех $t \geqslant 0$, поэтому неразложимости достаточно для того, чтобы за некоторое время $T$ установилась строгая положительность всех $p_{ij}(t)$. Это свойство выполняется в силу того, что время жизни в каждом состоянии распределено по закону $\text{Exp}(q_i)$, а условие неразложимости гарантирует наличие пути между любыми состояниями. Поскольку скачок в соседнее состояние возможен за сколь угодно малое время, из любого $i$ можно попасть в любое $j$ с ненулевой вероятностью за любое время $t > 0$. Далее доказательство сводится к дискретному случаю.
\end{tproof}

\begin{theorem}{Эргодическое свойство средних*}
    В конечной однородной неразложимой непрерывной цепи Маркова почти наверное выполняется:
    \begin{equation}
        \frac{1}{T} \int_0^T I(X(t)=j) \, dt \xrightarrow[T \to \infty]{\text{п.н.}} p_j,
    \end{equation}
    где $p_j$ — стационарная вероятность, выражающаяся через среднее время пребывания в состоянии $1/q_j$ и среднее время возврата $\mathbb{E}\tau_{jj}$:
    \begin{equation}
        p_j = \frac{1/q_j}{\mathbb{E}\tau_{jj}}.
    \end{equation}
\end{theorem}

\begin{note}
    Физический смысл $p_j = \frac{\mathbb{E} \text{Exp}(q_j)}{\mathbb{E} \tau_{jj}}$ заключается в том, что предельная доля времени нахождения в состоянии $j$ равна отношению среднего времени одного визита в это состояние к среднему времени между циклами возвращения.
\end{note}

\section{Примеры решения задач}

\begin{example}{Нахождение вероятностей состояний и переходной матрицы}
    Рассмотрим непрерывную цепь Маркова с двумя состояниями $S = \{0, 1\}$. Граф интенсивностей переходов задан следующим образом:
    \begin{center}
        \begin{tikzpicture}[>=stealth, node distance=3cm, on grid, auto]
            \node[state] (0) {0};
            \node[state] (1) [right=of 0] {1};
            \path[->] 
                (0) edge [bend left] node {$\lambda$} (1)
                (1) edge [bend left] node {$\mu$} (0);
        \end{tikzpicture}
    \end{center}
    Где $\lambda > 0$ и $\mu > 0$ --- заданные интенсивности. Требуется найти вектор вероятностей состояний $\pi(t)$ и матрицу переходных вероятностей $P(t)$.
\end{example}

\begin{eproof}
    \textbf{1. Формирование матрицы интенсивностей.}
    Согласно определению матрицы $Q$, внедиагональные элементы $q_{ij}$ равны интенсивностям соответствующих переходов. Диагональные элементы $q_{ii}$ определяются из условия равенства нулю суммы элементов каждой строки:
    \begin{equation}
        Q = \mat{-\lambda & \lambda \\ \mu & -\mu}.
    \end{equation}
    Транспонированная матрица, необходимая для уравнений Колмогорова, имеет вид:
    \begin{equation}
        Q^T = \mat{-\lambda & \mu \\ \lambda & -\mu}.
    \end{equation}

    \textbf{2. Решение системы дифференциальных уравнений.}
    Вектор вероятностей $\pi(t) = (\pi_0(t), \pi_1(t))^T$ удовлетворяет системе $\dot{\pi}(t) = Q^T \pi(t)$:
    \begin{equation}
        \begin{cases} 
            \dot{\pi}_0(t) = -\lambda \pi_0(t) + \mu \pi_1(t) \\ 
            \dot{\pi}_1(t) = \lambda \pi_0(t) - \mu \pi_1(t) 
        \end{cases}
    \end{equation}
    Воспользуемся условием нормировки $\pi_0(t) + \pi_1(t) = 1$, откуда $\pi_1(t) = 1 - \pi_0(t)$. Подставим это в первое уравнение:
    \begin{equation}
        \dot{\pi}_0(t) = -\lambda \pi_0(t) + \mu (1 - \pi_0(t)) = -(\lambda + \mu)\pi_0(t) + \mu.
    \end{equation}
    Это линейное неоднородное уравнение первого порядка. Общее решение соответствующего однородного уравнения: $\pi_{0, \text{одн}}(t) = C e^{-(\lambda + \mu)t}$. Частное решение ищем в виде константы (стационарное состояние $\dot{\pi}_0 = 0$):
    \begin{equation}
        0 = -(\lambda + \mu)\pi_0^* + \mu \implies \pi_0^* = \frac{\mu}{\lambda + \mu}.
    \end{equation}
    Тогда общее решение:
    \begin{equation}
        \pi_0(t) = \frac{\mu}{\lambda + \mu} + C e^{-(\lambda + \mu)t}.
    \end{equation}
    Из начального условия при $t=0$ находим $C = \pi_0(0) - \frac{\mu}{\lambda + \mu}$. Окончательно:
    \begin{equation}
        \pi_0(t) = \frac{\mu}{\lambda + \mu} + \left( \pi_0(0) - \frac{\mu}{\lambda + \mu} \right) e^{-(\lambda + \mu)t}, \quad \pi_1(t) = 1 - \pi_0(t).
    \end{equation}

    \textbf{3. Нахождение переходной матрицы $P(t)$.}
    Матрица $P(t)$ удовлетворяет уравнению $\dot{P}(t) = P(t)Q$ с начальным условием $P(0) = I$. Заметим, что каждая строка матрицы $P(t)$ — это вектор вероятностей состояний при специфическом начальном распределении (либо $(1, 0)$ или $(0, 1)$). 
    
    Для первой строки $(p_{00}(t), p_{01}(t))$ начальное условие $p_{00}(0) = 1$. Подставляя его в решение для $\pi_0(t)$:
    \begin{equation}
        p_{00}(t) = \frac{\mu}{\lambda + \mu} + \left( 1 - \frac{\mu}{\lambda + \mu} \right) e^{-(\lambda + \mu)t} = \frac{\mu}{\lambda + \mu} + \frac{\lambda}{\lambda + \mu} e^{-(\lambda + \mu)t}.
    \end{equation}
    Для второй строки $(p_{10}(t), p_{11}(t))$ начальное условие $p_{10}(0) = 0$:
    \begin{equation}
        p_{10}(t) = \frac{\mu}{\lambda + \mu} + \left( 0 - \frac{\mu}{\lambda + \mu} \right) e^{-(\lambda + \mu)t} = \frac{\mu}{\lambda + \mu} - \frac{\mu}{\lambda + \mu} e^{-(\lambda + \mu)t}.
    \end{equation}
    Остальные элементы находятся как $p_{01} = 1 - p_{00}$ и $p_{11} = 1 - p_{10}$.

    \textbf{4. Стационарный режим.}
    Поскольку цепь неразложима ($\lambda, \mu > 0$), существует единственное стационарное распределение $\pi$, которое является пределом $P(t)$ при $t \to \infty$. Решая алгебраическую систему $Q^T \pi = 0$:
    \begin{equation}
        \begin{cases} 
            -\lambda \pi_0 + \mu \pi_1 = 0 \\ 
            \pi_0 + \pi_1 = 1 
        \end{cases} \implies
        \pi_0 = \frac{\mu}{\lambda + \mu}, \quad \pi_1 = \frac{\lambda}{\lambda + \mu}.
    \end{equation}
    Этот результат полностью согласуется с пределом $t \to \infty$ в динамическом решении.
\end{eproof}

\begin{note}
    В данном примере мы явно увидели проявление эргодичности: независимо от начального состояния (выбора строки в $P(t)$ или начального вектора $\pi(0)$), система стремится к одному и тому же распределению вероятностей.
\end{note}

\begin{example}{Процесс рождения и гибели}
    Однородная непрерывная цепь Маркова называется процессом рождения — гибели, если её стохастический граф имеет линейную структуру, где переходы возможны только между соседними состояниями:
    \begin{center}
        \begin{tikzpicture}[>=Stealth, node distance=2.5cm, auto]
            % Состояния
            \node[state] (0) {0};
            \node[state] (1) [right=of 0] {1};
            \node[state] (2) [right=of 1] {2};
            \node[right=of 2] (dots) {$\dots$};
            \node[state] (N-1) [right=of dots] {$N-1$};

            % Стрелки рождения (вверх)
            \path[->] (0) edge [bend left] node {$\lambda_0$} (1);
            \path[->] (1) edge [bend left] node {$\lambda_1$} (2);
            \path[->] (2) edge [bend left] node {$\lambda_2$} (dots);
            \path[->] (dots) edge [bend left] node {$\lambda_{N-2}$} (N-1);

            % Стрелки гибели (вниз)
            \path[->] (1) edge [bend left] node {$\mu_1$} (0);
            \path[->] (2) edge [bend left] node {$\mu_2$} (1);
            \path[->] (dots) edge [bend left] node {$\mu_3$} (2);
            \path[->] (N-1) edge [bend left] node {$\mu_{N-1}$} (dots);
        \end{tikzpicture}
    \end{center}
    Здесь $\lambda_i > 0$ — интенсивности рождения, $\mu_i > 0$ — интенсивности гибели, а количество состояний $N$ может быть как конечным, так и бесконечным. Процесс является непрерывным аналогом случайных блужданий. Найдем для него явный вид стационарных вероятностей состояний $\pi_i$.
\end{example}

\begin{eproof}
    Прежде всего отметим, что стационарное распределение существует, является единственным и совпадает с предельным распределением вероятностей по эргодической теореме, поскольку рассматриваемая цепь образует замкнутый класс. Найдем стационарное распределение $\pi^\circ$. Матрица интенсивностей $Q$ для такого процесса является трехдиагональной:
    \begin{equation}
        Q = \begin{pmatrix}
        -\lambda_0 & \lambda_0 & 0 & 0 & \dots \\
        \mu_1 & -(\mu_1 + \lambda_1) & \lambda_1 & 0 & \dots \\
        0 & \mu_2 & -(\mu_2 + \lambda_2) & \lambda_2 & \dots \\
        \vdots & \vdots & \ddots & \ddots & \ddots
        \end{pmatrix}
    \end{equation}
    Для нахождения $\pi^\circ$ решим уравнение $\pi^{\circ T} Q = 0$. Распишем систему по столбцам:
    \begin{itemize}
        \item[1)] $-\pi_0 \lambda_0 + \pi_1 \mu_1 = 0 \implies \pi_1 = \pi_0 \frac{\lambda_0}{\mu_1}$
        \item[2)] $\pi_0 \lambda_0 - \pi_1 (\lambda_1 + \mu_1) + \pi_2 \mu_2 = 0$. Подставляя $\pi_0 \lambda_0 = \pi_1 \mu_1$, получаем:
        \begin{equation}
            \pi_1 \mu_1 - \pi_1 \lambda_1 - \pi_1 \mu_1 + \pi_2 \mu_2 = 0 \implies \pi_2 = \pi_1 \frac{\lambda_1}{\mu_2} = \pi_0 \frac{\lambda_0 \lambda_1}{\mu_1 \mu_2}
        \end{equation}
    \end{itemize}
    По индукции для любого $k \ge 1$ получаем общее уравнение баланса $\pi_i \lambda_i = \pi_{i+1} \mu_{i+1}$, откуда:
    \begin{equation}
        \pi_k = \pi_0 \prod_{i=1}^k \frac{\lambda_{i-1}}{\mu_i}
    \end{equation}
    Вероятность $\pi_0$ находится из условия нормировки $\sum_{k=0}^{N-1} \pi_k = 1$:

    \pagebreak

    \begin{equation}
        \pi_0 \left( 1 + \sum_{k=1}^{N-1} \prod_{i=1}^k \frac{\lambda_{i-1}}{\mu_i} \right) = 1 \implies \pi_0 = \left[ 1 + \sum_{k=1}^{N-1} \prod_{i=1}^k \frac{\lambda_{i-1}}{\mu_i} \right]^{-1}
    \end{equation}

\end{eproof}

\begin{definition}{Пуассоновский процесс (определение через процесс рождения)}
    Пусть $K(t), t \ge 0$ — непрерывная цепь Маркова с $S = \{0, 1, 2, \dots \}$, непрерывными справа траекториями, являющаяся процессом чистого рождения с постоянными интенсивностями:
    \begin{equation}
        \lambda_i = \lambda > 0, \quad \mu_i = 0 \quad \forall i.
    \end{equation}
    Тогда $K(t)$ называется пуассоновским процессом с интенсивностью $\lambda$.
\end{definition}

\begin{note}
    Поскольку пуассоновский процесс обладает независимыми приращениями, он автоматически удовлетворяет марковскому свойству. В данной интерпретации мы рассматриваем его как вырожденный случай процесса рождения-гибели, где «гибель» невозможна.
\end{note}

\section{Потоки событий}

\begin{definition}{Поток событий}
    \textit{Потоком событий} называется последовательность одинаково распределённых событий, происходящих одно за другим через случайные промежутки времени.
\end{definition}

\begin{center}
    \begin{tikzpicture}
        \draw[->, thick] (0,0) -- (10,0) node[right] {$t$};
        \node at (0,-0.3) {0};
        \draw[thick] (0,0.1) -- (0,-0.1);
        % Имитация случайных моментов наступления событий
        \foreach \x in {0.8, 1.5, 2.3, 2.5, 2.7, 3.8, 4.2, 4.8, 5.5, 6.7, 7.2, 8.4, 9.1}
            \draw[thick, MAIN] (\x, 0.15) -- (\x, -0.15);
    \end{tikzpicture}
\end{center}

Пусть $t_1, t_2 > 0$. Обозначим через $n(t_1, t_2)$ случайную величину, равную числу событий из потока, попавших в полуинтервал $[t_1, t_2)$.

\begin{definition}{Свойства потоков}
    \begin{enumerate}
        \item Поток называется \textbf{однородным (стационарным)}, если распределения случайных величин $n(t_1, t_2)$ и $n(t_1 + s, t_2 + s)$ совпадают для любого $s > 0$. Для таких потоков вводится обозначение $n(t) \doteq n(0, t)$.
        \item Рассмотрим моменты времени $0 = t_0 \le t_1 \le \dots \le t_{m+1}$ и соответствующие приращения $\xi_k = n(t_k, t_{k+1})$. Поток называется \textbf{потоком без последействия}, если совокупность случайных величин $\{\xi_0, \dots, \xi_m \}$ независима в совокупности.
    \end{enumerate}
\end{definition}



Дополнительно введем ограничение: между любыми двумя событиями в потоке проходит конечное (но случайное) время. Это означает, что после наступления события наступает интервал положительной длительности, в течение которого новых событий не происходит, но по истечении которого событие обязательно наступит. Это допущение эквивалентно требованию \textbf{непрерывности справа} траекторий процесса $n(t)$.

\subsection{Марковская природа потоков без последействия}

Исследуем структуру такого потока. Заметим, что:
\begin{enumerate}
    \item $n(t_k, t_{k+1}) = n(t_{k+1}) - n(t_k)$ независимы в совокупности по определению отсутствия последействия.
    \item Следовательно, $n(t)$ является процессом с \textbf{независимыми приращениями}.
    \item Любой процесс с независимыми приращениями обладает марковским свойством.
\end{enumerate}

Так как процесс $n(t)$ определен на непрерывном времени $t \in [0, +\infty)$, а его фазовое пространство дискретно ($S = \mathbb{N}_0$), то $n(t)$ представляет собой \textbf{непрерывную цепь Маркова}. Если поток однороден, то и цепь Маркова является однородной.

Тот факт, что траектории $n(t)$ непрерывны справа, делает эту цепь \textbf{стандартной}. Из этого вытекают фундаментальные аналитические свойства:
\begin{equation}
    P(t) \xrightarrow[t \to 0+]{} I, \quad P(t) \in C[0, +\infty).
\end{equation}

\subsection{Интенсивности переходов в потоке без последействия}

Рассмотрим вопрос об интенсивностях переходов в марковской цепи, описывающей поток событий. Исходя из свойств траекторий процесса $n(t)$, мы можем сделать ряд важных выводов о структуре матрицы интенсивностей.

\begin{statement}
Для марковской цепи, соответствующей однородному потоку без последействия, интенсивности переходов $\lambda_{ij}$ не зависят от текущего состояния $i$ и отличны от нуля только для переходов в непосредственно следующее состояние ($j = i+1$). При этом диагональные элементы $\lambda_{ii}$ конечны и равны по модулю интенсивности потока.
\end{statement}

\begin{sproof}
1. \textbf{Анализ переходных вероятностей.} 
Так как траектории процесса $n(t)$ являются неубывающими по определению (количество событий не может уменьшаться), то переходные вероятности $p_{ij}(t)$ отличны от нуля только при $j \ge i$. Для стандартной цепи Маркова при $j > i$ существуют конечные интенсивности переходов:
$$ \exists \lambda_{ij} = \dot{p}_{ij}(0), \quad 0 \le \lambda_{ij} < \infty. $$
Для диагональных элементов $\lambda_{ii} = \dot{p}_{ii}(0)$ априори допускается значение $-\infty$, однако покажем, что в данном случае они конечны.

2. \textbf{Строгая положительность интенсивности выхода.}
Заметим, что $\lambda_{ii} < 0$ (так как $p_{ii}(t)$ убывает от 1). Предположим, что $\lambda_{ii} = 0$. Тогда, воспользовавшись мультипликативным свойством и перейдя к пределу, получим:
$$ p_{ii}(t) = \left[ p_{ii}\left(\frac{t}{n}\right) \right]^n = \left( 1 + \dot{p}_{ii}(0) \frac{t}{n} + o\left(\frac{t}{n}\right) \right)^n = \left( 1 + 0 \cdot \frac{t}{n} + o\left(\frac{t}{n}\right) \right)^n \xrightarrow[n \to \infty]{} 1. $$
Следовательно, $p_{ii}(t) \equiv 1$ для всех $t \ge 0$. Это означает, что система никогда не покидает состояние $i$, что прямо противоречит условию о том, что между событиями проходит конечное время. Таким образом, $\lambda_{ii} < 0$.

3. \textbf{Независимость от состояния (однородность).}
Предположим, что $\lambda_{ii}$ конечна. Рассмотрим вероятность того, что на интервале длины $h$ не произойдет ни одного события:
$$ p_{ii}(h) = \mathbb{P}(n(t+h) = i \mid n(t) = i) = \mathbb{P}(n(t+h) - n(t) = 0 \mid n(t) = i). $$
В силу свойства независимости приращений, данная вероятность совпадает с безусловной вероятностью того, что за время $h$ приращение процесса составит $0$:
$$ p_{ii}(h) = \mathbb{P}(n(h) = 0). $$
Отсюда следует, что величина $p_{ii}(h)$ (а значит, и её производная в нуле $\lambda_{ii}$) не зависит от индекса состояния $i$. Обозначим $-\lambda_{ii} \doteq \lambda$. Как было показано ранее, время пребывания в состоянии в таком случае распределено экс­по­нен­ци­аль­но: $\tau_i \sim \text{Exp}(\lambda)$.

4. \textbf{Отсутствие групповых скачков.}
Рассмотрим переход из $i$ в $i+1$. Аналогично предыдущему пункту:
$$ p_{i, i+1}(h) = \mathbb{P}(n(t+h) - n(t) = 1 \mid n(t) = i) = \mathbb{P}(n(h) = 1). $$
Это означает, что интенсивность перехода $\lambda_{i, i+1}$ также не зависит от $i$. Обозначим её за $\lambda$. Тогда:
$$ \frac{p_{i, i+1}(h)}{h} \xrightarrow[h \to 0]{} \lambda. $$
Рассмотрим переходы типа $i \to i+k$ при $k > 1$ (например, $0 \to 2$). По условию, между любыми двумя событиями проходит конечное время. Это исключает возможность мгновенного скачка через несколько состояний (групповое поступление событий). Следовательно, такие интенсивности равны нулю: $\lambda_{ij} = 0$ при $j > i+1$.
\end{sproof}

Визуально структуру переходов в такой цепи можно представить следующей диаграммой:

\begin{figure}[h!]
    \centering
    \begin{tikzpicture}[
        scale=1.5,
        >={Stealth[length=3mm]},
        state/.style={circle, draw, minimum size=0.8cm, thick},
        every node/.style={font=\small}
    ]
        % Состояния
        \node[state] (s0) at (0,0) {0};
        \node[state] (s1) at (2,0) {1};
        \node[state] (s2) at (4,0) {2};
        \node (dots) at (5.5,0) {\Large $\dots$};

        % Разрешенные переходы (интенсивности в кружочках над стрелками)
        \draw[->, thick] (s0) to [bend left=40] node[midway, circle, draw, inner sep=1.5pt, yshift=0.3cm] {$\lambda$} (s1);
        \draw[->, thick] (s1) to [bend left=40] node[midway, circle, draw, inner sep=1.5pt, yshift=0.3cm] {$\lambda$} (s2);
        \draw[->, thick] (s2) to [bend left=40] node[midway, circle, draw, inner sep=1.5pt, yshift=0.3cm] {$\lambda$} (dots);

        % Запрещенный переход (пунктир снизу)
        \path[draw, ->, dashed, thick, gray] (s0) to [bend right=45] coordinate[midway] (cross) (s2);
        
        % Классический красный крестик поверх линии
        \draw[red, ultra thick] ($(cross)+(-0.15, -0.15)$) -- ($(cross)+(0.15, 0.15)$);
        \draw[red, ultra thick] ($(cross)+(-0.15, 0.15)$) -- ($(cross)+(0.15, -0.15)$);

    \end{tikzpicture}
    \caption{Схема интенсивностей переходов пуассоновского процесса.}
\end{figure}

\begin{note}
Таким образом, мы получили непрерывную цепь Маркова со счетным фазовым пространством $S = \mathbb{N}_0$, стартующую из нуля, имеющую независимые приращения и экспоненциальные времена между скачками. Данная конструкция (схема переходов $0 \to 1 \to 2 \to \dots$) полностью определяет \textbf{пуассоновский процесс}.
\end{note}

\subsection{Ординарность и эквивалентные определения пуассоновского процесса}

В предыдущих рассуждениях мы существенно опирались на допущение о том, что между любыми двумя последовательными событиями в потоке проходит конечное время, и интенсивности выхода из состояний удовлетворяют условию $0 < \lambda_i < \infty$. Данные условия гарантируют \textit{стандартность} цепи Маркова, а также определяют структуру времени пребывания в состояниях и вероятности скачков.

Если мы рассматриваем поток, в котором на любом конечном интервале времени происходит конечное число событий, условие $\lambda_i < \infty$ выполняется автоматически. При этом требование конечности межсобытийных интервалов эквивалентно требованию \textit{непрерывности справа} траекторий процесса $n(t)$.

Формально условия на структуру приращений процесса можно переписать через свойство ординарности.

\begin{statement}
Для однородного потока событий с интенсивностью $0 < \lambda < \infty$ вероятности изменения состояния на малом интервале времени $h \to 0$ представимы в виде:
\begin{equation}
    \begin{cases}
        \mathbb{P}(n(t, t+h) = 1) = \lambda h + o(h), \\
        \mathbb{P}(n(t, t+h) = 0) = 1 - \lambda h + o(h).
    \end{cases}
\end{equation}
Отсюда, в силу нормировки, следует, что вероятность наступления двух и более событий за малое время пренебрежимо мала:
$$ \mathbb{P}(n(t, t+h) \ge 2) = o(h), \quad h \to 0. $$
Данное свойство называется \textbf{ординарностью} потока и исключает возможность группового поступления событий.
\end{statement}

Теперь мы можем дать строгое определение рассматриваемому процессу.

\begin{definition}{Пуассоновский процесс}
Пусть дан однородный ординарный поток без последействия с интенсивностью $\lambda$. Тогда случайный процесс $K(t) \doteq n(0, t)$, равный числу событий, произошедших на полуинтервале $[0, t)$, называется \textbf{пуассоновским процессом} с интенсивностью $\lambda$ (или простейшим пуассоновским потоком).
\end{definition}

\begin{note}
В случае отказа от условия однородности интенсивность становится функцией времени:
$$ \lambda(t) \doteq \lim_{h \to 0} \frac{\mathbb{P}(n(t, t+h) > 0)}{h}. $$
Тогда свойство ординарности принимает локальный вид:
\begin{align*}
    \mathbb{P}(n(t, t+h) = 1) &= \lambda(t)h + o(h), \\
    \mathbb{P}(n(t, t+h) = 0) &= 1 - \lambda(t)h + o(h).
\end{align*}
Такой процесс называется \textit{неоднородным пуассоновским процессом}.
\end{note}

Подводя итог изучению данного класса процессов, перечислим четыре эквивалентных подхода к его определению, которые мы встретили в рамках курса:

\begin{enumerate}
    \item \textbf{Аксиоматическое определение:} через свойства независимых и стационарных приращений, распределенных по закону Пуассона.
    \item \textbf{Явная конструкция:} как процесс, в котором времена между скачками независимы и распределены экспоненциально с параметром $\lambda$.
    \item \textbf{Процесс гибели и размножения:} как частный случай (чистый процесс рождения) марковской цепи с непрерывным временем и линейной структурой переходов.
    \item \textbf{Потоковое определение:} как однородный ординарный поток без последействия.
\end{enumerate}